\section*{Experimental Design, Materials and Methods}

\subsection*{Scope and design}
PhishVN-SMS is a \emph{pilot}: \FILL{30} contributors at around \FILL{100} messages each,
about three thousand messages, enough to study labelling, templates and splitting, and not
enough to describe Vietnamese smishing. The target mix (roughly 800 phishing, 450 spam, 1{,}800
legitimate) is a designed 1:2 malicious-to-legitimate ratio, not the prevalence of any inbox;
all reported metrics are per class, and nothing depends on this base rate. Messages are what is
collected, but templates are what the evaluation counts: the figure that decides what the corpus
can say is the number of distinct malicious templates (floor: 100), because another copy of a
template already held adds nothing a template-disjoint split can use. Collection watched the
template count as it went and stopped adding a contributor's near-duplicates before stopping
their novel messages.

Two kinds of message were collected, and only two. \emph{Scam and spam} of any sender,
broadcast to strangers by someone with no privacy interest in them; and \emph{legitimate
messages sent by a machine} (one-time codes, bank and telecommunications notifications, school
and application notices) from a brandname or a service shortcode. \emph{Person-to-person
messages were never collected.} A message a friend or relative sent is theirs, and they are not
present to consent; a contributor can give away what is their own and nobody can consent on a
sender's behalf. The rule contributors were given is mechanical: a message from a personal phone
number is not sent, whatever it says, even when it is plainly a scam.

\subsection*{Consent design}
Contributors may be students, and the lead researcher teaches. A student asked by a lecturer
for phone data is not free to refuse, and a consent form does not fix that by itself. The
design removes the asymmetry instead of documenting it:
\begin{enumerate}[leftmargin=1.6em,itemsep=1pt,topsep=2pt]
  \item The invitation came from a colleague with no assessment role over the cohort.
  \item Nothing was attached to participation: no marks, credit, bonus or attendance, and no
  mention in class of who took part.
  \item Submission was anonymous. The researchers never learn who participated, so declining
  costs nothing and cannot be noticed.
  \item Withdrawal needed no reason and was possible until deposit. After deposit the
  researchers cannot identify which messages to remove, and the consent form says so in those
  words instead of promising a deletion that cannot be performed.
\end{enumerate}
The consent form (\texttt{CONSENT\_vi.md}, in Vietnamese; the English translation is for
reviewers and the Vietnamese governs) states in plain language what is asked for and what is
not, that publication is permanent, that the screenshot route exposes the image to the
researcher, and the principal risk: a personal detail the filter missed and the contributor did
not notice on review. It asks for no signature, since a signature would turn an anonymous
submission into an identified one.

\subsection*{Collection and redaction}
Figure~\ref{fig:pipeline} shows the pipeline from the contributor's phone to the two published
files. Every stage reads files and writes files, so the whole of it can be rerun and checked.

\begin{figure}[t]
  \centering
  \includegraphics[width=\linewidth]{pipeline.pdf}
  \caption{PhishVN-SMS pipeline. On the contributor's device, messages are pasted into an
  offline redaction page (the raw text never leaves the phone) or sent as cropped, blacked-out
  screenshots. On the author's machine, screenshots are OCR-transcribed, hand-checked, redacted
  by the same rules and deleted; rows are validated and ingested; two annotators label blind
  and disagreements are adjudicated; templates are computed and the split assigned per
  template; the publishing script projects the schema's columns and checks its rules. The
  external benchmark enters by a separate import, is held out as \texttt{split = external},
  and ships as its own file with the source's labels.}
  \label{fig:pipeline}
\end{figure}

\paragraph{Primary route: paste.} The contributor opens a single offline HTML page
(\texttt{redact.html}) in the phone's own browser. The page makes no network request and works
in airplane mode. Each message is copied from the messaging app and pasted in, so the text
arrives character for character with no screenshot or OCR involved. The page replaces personal
data with the placeholders listed under \emph{Data Description}; the contributor reviews every
message after redaction and before sending, and drops any they are unsure about. The raw text
never leaves the phone.

\paragraph{What is kept, and what is masked.} Links are kept. A link was broadcast to strangers
and is not personal data; only personal data \emph{inside} a link is masked: a phone number, an
e-mail address, or a name or account in a query string (\texttt{?sdt=}, \texttt{?ten=},
\texttt{?stk=}). A path segment that identifies the recipient without saying so cannot be
recognised by rule, so contributors were asked to check each link and drop a message whose link
carries their own details. Broadcast amounts are kept: a prize, a fee or a price in a scam or an
advert is a signal. The contributor's own money is masked: a balance, a signed account movement
(\texttt{+1.500.000VND}) and an amount beside a transaction marker are sensitive personal data
and become \texttt{<AMOUNT>}. Normalisation stops there; spelling, punctuation and Unicode are
untouched.

\paragraph{Alternative route: screenshots.} A message that cannot be copied may be sent as a
screenshot, cropped to the sender and the message, with the contributor's own details blacked
out on the image, in one \texttt{.zip} file. On this route the researcher does see an unredacted
image, and the consent form says so. On receipt, images are re-encoded from pixels alone so that
no file name or metadata survives, transcribed by OCR (\FILL{engine and version}), checked by
hand against each image, redacted by the same rules as the paste route (the transcription script
reads the rules out of \texttt{redact.html}, so the two routes cannot drift), and deleted once
checked. \texttt{capture = screenshot} marks every such row.

\paragraph{Ingest.} Each submission row is validated (reviewed flag, sender type, month format,
capture route) before being appended to the submissions file. The collector refuses to run until
the protocol carries no draft status and the consent form has no unfilled blank; a unit test
asserts that it does.

\subsection*{Label scheme}
Three classes, defined for annotators as in Table~\ref{tab:labels}.

\begin{table}[h]
\centering\small
\caption{The label scheme.}
\label{tab:labels}
\begin{tabular}{@{}l p{11.5cm}@{}}
\toprule
Label & Definition \\
\midrule
\texttt{legitimate} & a genuine message from the organisation it appears to come from \\
\texttt{spam} & unsolicited or promotional, with no evidence of intent to obtain credentials, money or information by deception \\
\texttt{phishing} & impersonates an organisation or person, or seeks credentials, a one-time code, a payment, a call-back or an app install under a false pretext \\
\bottomrule
\end{tabular}
\end{table}

\emph{A message is not made \texttt{phishing} because it looks suspicious.} Evidence of
deceptive intent is required; where it is absent the message is \texttt{spam}.
\texttt{uncertain} is a fourth value an annotator may use and the published
\texttt{final\_label} never carries: it sends a message to adjudication. For a binary experiment
the published collapse rule is \texttt{legitimate} $+$ \texttt{spam} $\to 0$, \texttt{phishing}
$\to 1$.

\subsection*{Annotation}
Two human annotators labelled every contributed message independently. The labelling tool shows
an annotator the text and nothing else: not the other annotator's label, not the contributor's
own label, not the sender or any metadata. Every disagreement and every \texttt{uncertain} went
to an adjudicator, whose decision is \texttt{final\_label} and is recorded with a name and a
one-line note. Agreements were finalised without an adjudicator.

Three things are reported whatever they show: Cohen's $\kappa$~\cite{cohen_sms} between the two
annotators before adjudication (\FILL{}), the disagreement rate per class pair (\FILL{}), and
the share adjudicated (\FILL{}). A low $\kappa$ is a result about how separable these classes
are in practice; it is published as one and not repaired by relabelling until the annotators
agree. Both original labels ship unmodified so the agreement can be recomputed.
\FILL{annotator background and training; whether the adjudicator was one of the annotators or a
third person. The working file carries a machine pre-annotation column that the labelling tool
never shows to an annotator and that is excluded from $\kappa$; state this.}

\subsection*{Templates}
Grouping is computed, not assigned by hand. It runs on a \emph{normalized view} of the text
that exists only for this step: lower-cased, whitespace collapsed, placeholders kept, and every
URL and money amount replaced by a \texttt{<URL>} or \texttt{<AMT>} token. The published text
keeps both, but for grouping they are the slots a campaign rotates, and leaving them in place
splits one template into many. The amount patterns for this step extend the redaction rules
with the short forms (one to three digits plus a thousand/million/billion suffix) that redaction
rightly ignores, since two digits identify nobody, but grouping must not.

Each normalized message becomes the set of its word 4-shingles~\cite{broder_sms}. Pairs at or
above Jaccard similarity $\tau = 0.8$ are linked, and a template is a connected component of
that graph. Results at $\tau = 0.7$ and $\tau = 0.9$ are reported beside the primary so that
chaining is visible. The shingling, thresholds and connected components are the method the
sibling kit-reuse study registers~\cite{kitreuse_sms}; the normalization is this protocol's one
departure from it, adopted after a pilot batch showed the failure it repairs: a scam sent twice
with one edited character in the domain, or a different lure amount, fell below $\tau$ even at
0.7, because in a message of eight words one changed word destroys four shingles. Counted as
different templates, the two copies then land on opposite sides of a template-disjoint split,
which is the leakage the split exists to prevent.

\subsection*{Splits}
\emph{Every message of one template lies in one split.} The published corpus this project
audits~\cite{qavn_sms} does not do this: 47 of its 597 test rows (7.9\%) repeat a training
text. A fixed 70/15/15 template-disjoint split ships with the corpus. It is assigned per
template, never per message, by a greedy search over 200 seeded restarts (seed 1602), scored on
split size and label mix; the seed and restart count are constants in the script, so the split
reproduces from the working file alone, and a unit test asserts that no template crosses a
split.

Three evaluations are specified for the paper that uses the corpus, in decreasing order of what
they can support:
\begin{enumerate}[leftmargin=1.6em,itemsep=1pt,topsep=2pt]
  \item \emph{Cross-validation over template groups} is the primary evaluation. The fixed
  split is for presentation; with a pilot this size its test slice holds too few phishing
  messages to carry a conclusion, and the paper says so where it reports it.
  \item \emph{Leave-one-contributor-out}: train on all contributors but one, test on the one.
  With \FILL{n} contributors that is as many folds, each a single person, so the folds are
  reported individually and never averaged. It answers ``does this transfer to an inbox it has
  not seen'' qualitatively and estimates nothing. It is computed from the private
  \texttt{participant\_id} and reported without the mapping.
  \item \emph{The external benchmark}, below.
\end{enumerate}

\subsection*{The external benchmark}
The Quality-Assured Vietnamese SMS Phishing Dataset~\cite{qavn_sms} (CC BY 4.0; 2{,}991
messages; binary labels, 2{,}193 \texttt{benign} and 798 \texttt{scam}) is held out entirely:
every row is \texttt{split = external}, never train, validation or test. The copy used is
pinned by the SHA-256 of its \texttt{full\_dataset.csv} (Specifications Table); a copy
reproduces this benchmark only if its hash matches. No repository revision was recorded at
download (2026-08), so the content hash is the anchor.

Its placeholder tokens were mapped to this corpus's by an allowlist in the import script. Rows
that still failed the redaction rule after mapping were dropped and counted, never repaired:
310 of 2{,}991, most carrying an unmasked one-time code in a bank message. The loaded benchmark
is therefore 2{,}681 rows (1{,}910 \texttt{benign}, 771 \texttt{scam}). Its labels are its own;
this project does not re-annotate them and does not treat its three-class scheme as a
correction of the source's two-class one.

A model trained on \texttt{sms\_dataset.csv} is scored on the benchmark after collapsing its
prediction to the source's two classes (\texttt{legitimate} $\to$ \texttt{benign};
\texttt{spam}, \texttt{phishing} $\to$ \texttt{scam}), against the source's own labels. The
paper also reports, once, \emph{how far the two label schemes diverge}. The source's boundary
separates \emph{genuine sender} from \emph{scam}, so its \texttt{benign} class carries the
operators' own promotional messages, which this corpus's scheme calls \texttt{spam} whoever
sent them. In a model-assisted pass confirmed by a hand audit of a sample, \FILL{about 50\%} of
the source's \texttt{benign} rows read as \texttt{spam} under the three-class scheme, while
\FILL{about 0\%} of its \texttt{scam} rows read as \texttt{legitimate}. That divergence is the
reason the two corpora are kept in separate files rather than merged: a single pooled label
column would bury it. A template-grouping check confirms that no external template coincides
with a training template, so the benchmark is genuinely unseen.

\subsection*{Publishing}
The two published files are a mechanical projection of the private working file: the
publishing script keeps exactly the schema's columns, derives the four \texttt{has\_*} flags
from the patterns in \texttt{redact.html}, and refuses to ship a train/validation/test row with
an \texttt{uncertain} or empty \texttt{final\_label}, a missing \texttt{template\_id} or a
missing \texttt{split}. A test suite holds the redaction routes to one list of cases and the
later stages to their invariants: rotated slots group as one template, no template crosses a
split, agreements finalise without an adjudicator, the projection ships exactly the schema's
columns. The pipeline uses the Python~3 standard library only, apart from the OCR engine of the
screenshot route; all scripts, the redaction page and the tests are in the repository at
release \FILL{tag}.
